% TOP_PROBLEM: {"id":30007571,"problem_number":"LOCAL-30007571","title":"Reduced-game validity","category_id":21,"category":{"id":21,"name":"economics","display_name":"Economics","description":"Open economic theory statements, published research agendas, and proposed scoped specifications; question provenance and historical priority are qualified per record.","slug":"economics","order_index":21,"created_at":"2026-10-08T00:00:00Z"},"status":"research_question","external_url":null,"legacy_ids":[],"published":false,"statement":"Determine how to allocate computational effort between restricted-game solving, strategy expansion, and deviation validation to minimize certified full-game exploitability.","economics":{"original_id":60,"field":"Game theory and strategic interaction","kind":"design","formalization_basis":"adapted_research_question","source_status":"research_agenda","priority_status":"earliest_found","priority_note":"This is the earliest source in which the relevant agenda was verified during this audit; absolute priority has not been established.","earliest_verified_reference":{"authors":["Michael P. Wellman","Karl Tuyls","Amy Greenwald"],"title":"Empirical Game-Theoretic Analysis: A Survey","year":2025,"url":"https://arxiv.org/pdf/2403.04018","doi":null,"locator":"§6 Statistical Reasoning in Empirical Games, pp.1045–1052; §7.3 Frontiers in Strategy Exploration, pp.1057–1058","evidence_summary":"The survey discusses adaptive statistical conclusions and identifies unresolved theoretical characterization of strategy-exploration solvers."},"current_reference":{"authors":["Michael P. Wellman","Karl Tuyls","Amy Greenwald"],"title":"Empirical Game-Theoretic Analysis: A Survey","year":2025,"url":"https://arxiv.org/pdf/2403.04018","doi":null,"locator":"§6 Statistical Reasoning in Empirical Games, pp.1045–1052; §7.3 Frontiers in Strategy Exploration, pp.1057–1058","evidence_summary":"The survey discusses adaptive statistical conclusions and identifies unresolved theoretical characterization of strategy-exploration solvers."},"formalization_note":"The budget-allocation problem is adapted from §7.3. Certification requires an explicitly justified best-response or coverage guarantee; no implication from restricted equilibrium alone is claimed."},"statusQualification":"Published question or research agenda verified at the cited locator. The mathematical specification is adapted for this catalogue and includes additional assumptions; source publication does not establish first-ever priority or certify this exact model remains unsolved. The budget-allocation problem is adapted from §7.3. Certification requires an explicitly justified best-response or coverage guarantee; no implication from restricted equilibrium alone is claimed.","statusReviewedAt":"2026-10-08"}
% FIRST_POSTED: 2026-10-08
% ENTRY_KIND: statement_only
% !TeX program = lualatex
\documentclass[11pt]{article}
\usepackage{fontspec}
\usepackage[a4paper,margin=25mm]{geometry}
\usepackage{amsmath,amssymb,mathtools}
\usepackage{xurl}
\usepackage[hidelinks]{hyperref}
\newcommand{\sref}[2]{\hyperlink{source:#1:#2}{[S#2]}}
\setlength{\emergencystretch}{3em}
\begin{document}
% problemId:30007571
\section*{Reduced-game validity}\label{30007571}
\subsection{Definitions and mathematical statement}
Let a normal-form game have compact strategy sets \(S_i\), bounded payoffs \(u_i\), and a simulator for expected utility. A restricted game uses finite subsets \(X_i\subseteq S_i\), and a candidate product distribution \(\sigma\) is supported on \(X=\prod_iX_i\). Define restricted exploitability \(e_X(\sigma)\) by maximizing deviations only over \(X_i\), and full exploitability by
\[
 e_S(\sigma)=\max_i\sup_{s_i\in S_i}
 [u_i(s_i,\sigma_{-i})-u_i(\sigma)].
\]
An approximate best-response oracle returns a candidate deviation and a statistically valid upper bound on the unobserved supremum; state its optimization error and probability of failure. Determine how to allocate a fixed total computation budget between enlarging \(X_i\), solving the restricted game, simulating payoffs, and validating deviations so that the final upper confidence bound on \(e_S\) is smallest. Identify conditions under which successive expansion provably improves full-game regret, and quantify the loss caused by approximation or restricted oracle access. Without an oracle guarantee or explicit regularity and coverage assumptions, low \(e_X\) alone cannot certify low \(e_S\); any positive result must include the missing assumption. Evaluate proposed solver policies on held-out games under a common validation protocol. This operationalizes the source’s unresolved strategy-exploration characterization, while retaining established restricted-versus-full regret relationships.

\par\textbf{Formulation and scope.} The budget-allocation problem is adapted from §7.3. Certification requires an explicitly justified best-response or coverage guarantee; no implication from restricted equilibrium alone is claimed.

\par\textbf{Open-question provenance.} An explicit question or research agenda was verified at the source locator. This does not by itself certify that every later version remains unresolved \sref{30007571}{1}.

\par\textbf{Historical priority.} This is the earliest source in which the relevant agenda was verified during this audit; absolute priority has not been established.

\subsection{Short English statement}
Determine how to allocate computational effort between restricted-game solving, strategy expansion, and deviation validation to minimize certified full-game exploitability.

\subsection{Sources}
\begin{itemize}\raggedright
\item[\textnormal{[S1]}]\hypertarget{source:30007571:1}{} Michael P. Wellman, Karl Tuyls, Amy Greenwald. 2025. Empirical Game-Theoretic Analysis: A Survey. §6 Statistical Reasoning in Empirical Games, pp.1045–1052; §7.3 Frontiers in Strategy Exploration, pp.1057–1058.
\url{https://arxiv.org/pdf/2403.04018}.
{\small\textit{Earliest verified question/agenda reference; historical priority qualified below. The survey discusses adaptive statistical conclusions and identifies unresolved theoretical characterization of strategy-exploration solvers.}}
\end{itemize}
\end{document}
